\documentclass[12pt]{article}
\usepackage[utf8]{inputenc}
\usepackage[T1]{fontenc}
\usepackage{hyperref}
\title{Yawa Language Resources for K-12 Curriculum Development}
\author{Compiled by Claude Sonnet 5.5 (claude-sonnet-5-5) for Andr\'as Kornai}
\date{}
\begin{document}
\maketitle

\section{Introduction}
Yawa (ISO 639-3 yva) is spoken on Yapen Island, Papua, Indonesia, and belongs to the small Yawa-Saweru family. Wikipedia gives about 6,000 speakers for 1987. The reference description gives about 10,000. This report lists only what appeared in search results.

\section{Primary Educational Resources}
No primers, textbooks or school programmes were found in the searches.

\section{Linguistic and Reference Materials}
\begin{itemize}
\item Jones et al. (1989), a Yawa vocabulary (\emph{Perbendaharaan kata bahasa Yawa}), published in Jayapura by Cenderawasih University and SIL, as listed in ASJP: \url{https://asjp.clld.org/languages/YAWA}
\item Jones (1986), ``The Dialects of Yawa'', Papers in New Guinea Linguistics 25, Glottolog reference: \url{https://glottolog.org/resource/reference/id/7728}
\item PhD thesis \emph{The Papuan languages of Yapen Island: history and typology}, covering the phonology, morphology and syntax of Yawa-Saweru varieties: \url{https://research.vu.nl/en/publications/the-papuan-languages-of-yapen-island-history-and-typology/}
\item Overview of the Yawa languages: \url{https://en.wikipedia.org/wiki/Yawa_languages}
\item Yawa numeral system data: \url{https://lingweb.eva.mpg.de/channumerals/Yawa.htm}
\end{itemize}

\section{Religious and Community Resources}
None found in the searches.

\section{Digital and Audio Resources}
None found in the searches.

\section{Curriculum Development Context}
The reference description says Yawa is taught as a subject in some schools. The searches did not confirm this or surface any materials for it. A curriculum would start from the vocabulary and the thesis above and from local community work.

\end{document}